\section{Thiết lập thực nghiệm}
\label{sec:experimental-setup}

\subsection{Mục tiêu đánh giá}

Thực nghiệm được thiết kế để đánh giá toàn bộ pipeline từ retrieval đến câu
trả lời cuối cùng, đồng thời tách riêng chất lượng của từng thành phần. Năm
câu hỏi nghiên cứu được đặt ra như sau:

\begin{description}
    \item[RQ1 --- Chất lượng retrieval.]
    Hệ thống có truy hồi đúng rule, event và article liên quan đến câu hỏi
    hay không?

    \item[RQ2 --- Chất lượng causal path.]
    Hệ thống có khôi phục đúng thứ tự các event và rule trên gold causal
    path hay không? Nếu gold path đã xuất hiện trong tập ứng viên, hệ thống
    có xếp nó ở vị trí đầu tiên hay không?

    \item[RQ3 --- Chất lượng xác minh.]
    Query-aware verifier có phân loại đúng các claim được hỗ trợ và các claim
    về cạnh trực tiếp cần bị bác bỏ hay không?

    \item[RQ4 --- Chất lượng câu trả lời.]
    Answer generator có trả đúng loại thông tin được hỏi, giữ được nội dung
    của gold answer và trích dẫn đúng điều luật hay không?

    \item[RQ5 --- Ảnh hưởng của verification engine.]
    Structural SCM và node-deletion path ablation khác nhau như thế nào về
    chất lượng đầu ra, thời gian xử lý và mức độ chi tiết của trace?
\end{description}

Ngoài kết quả tổng thể, nghiên cứu phân tích metric theo question type, độ
khó và yêu cầu counterfactual để xác định các nhóm lỗi chính.

\subsection{Giao thức đánh giá}

Thực nghiệm sử dụng toàn bộ benchmark gồm \(250\) câu hỏi. Không có bước
huấn luyện hoặc fine-tuning trên benchmark. Hệ thống chỉ sử dụng câu hỏi làm
đầu vào tại thời điểm inference; các trường evaluation được giữ riêng cho
Step 7.

Không thực hiện chia train--validation--test vì các thành phần chính của
pipeline không được huấn luyện trên benchmark. Tuy nhiên, benchmark được sinh
từ cùng graph mà retriever sử dụng, nên kết quả được xem là đánh giá
closed-corpus chứ không phải đánh giá khả năng tổng quát hóa ngoài miền.

Mỗi cấu hình chạy lần lượt các bước:

\begin{enumerate}
    \item truy hồi event, rule và causal path;
    \item xác minh path và claim;
    \item sinh câu trả lời extractive;
    \item chuẩn hóa prediction về schema chung;
    \item tính metric bằng Step 7.
\end{enumerate}

Evaluator được chạy với tùy chọn \texttt{strict}. Do đó, quá trình đánh giá
sẽ dừng nếu thiếu prediction hoặc nếu prediction chứa ID không thuộc
benchmark. Cả hai paired run đều phải có đủ \(250\) prediction và không có
pipeline error mới được đưa vào so sánh.

\subsection{Phiên bản các thành phần}

Các phiên bản được sử dụng trong thực nghiệm chính được trình bày trong
Bảng~\ref{tab:versions}.

\begin{table}[t]
    \centering
    \caption{Phiên bản các thành phần trong pipeline thực nghiệm.}
    \label{tab:versions}
    \begin{tabular}{ll}
        \hline
        \textbf{Thành phần} & \textbf{Phiên bản} \\
        \hline
        Benchmark & \texttt{2.0} \\
        Step 4 verifier & \texttt{4.1-query-aware-legal-scm} \\
        Step 5 answer generator & \texttt{5.1-query-aware-grounded-answer} \\
        Batch runner & \texttt{2.1-query-aware-answer} \\
        Step 7 evaluator & \texttt{2.0-current-pipeline-schema} \\
        \hline
    \end{tabular}
\end{table}

Phiên bản được lưu trực tiếp trong metadata của artifact. Prediction không
có đúng version provenance không được xem là thuộc paired experiment chính.

\subsection{Môi trường thực thi}

Hai paired run được thực hiện trong cùng một môi trường server và sử dụng
cùng bộ resource. Để bảo đảm khả năng tái lập, phiên bản cuối của paper cần
bổ sung đầy đủ thông tin phần cứng và phần mềm trong
Bảng~\ref{tab:environment}.

\begin{table}[t]
    \centering
    \caption{Môi trường thực thi. Các trường phần cứng cần được điền từ
    server đã chạy thí nghiệm trước khi công bố.}
    \label{tab:environment}
    \begin{tabular}{ll}
        \hline
        \textbf{Thành phần} & \textbf{Thông tin} \\
        \hline
        Hệ điều hành & \textit{[Bổ sung từ server]} \\
        Phiên bản Python & \textit{[Bổ sung từ server]} \\
        CPU & \textit{[Bổ sung từ server]} \\
        RAM & \textit{[Bổ sung từ server]} \\
        GPU, nếu sử dụng & \textit{[Bổ sung từ server]} \\
        Embedding model & \texttt{BAAI/bge-m3} \\
        Vector index & FAISS \\
        Answer provider & Extractive \\
        \hline
    \end{tabular}
\end{table}

BGE-M3 được sử dụng cho semantic retrieval và semantic mapping
\cite{chen2024bgem3}. FAISS được sử dụng để tìm kiếm lân cận gần trên causal
memory \cite{johnson2019billion}.

Thông tin CPU, RAM, GPU và phiên bản thư viện không được suy đoán từ artifact.
Các trường này cần được lấy trực tiếp từ server thực nghiệm. Điều này đặc
biệt quan trọng khi báo cáo runtime.

\subsection{Cấu hình retrieval}

Hai paired run sử dụng cùng embedding, graph, causal memory và FAISS index.
Các hyperparameter của Step 3 được trình bày trong
Bảng~\ref{tab:retrieval-hyperparameters}.

\begin{table}[t]
    \centering
    \caption{Hyperparameter của multi-hop causal retriever.}
    \label{tab:retrieval-hyperparameters}
    \begin{tabular}{lr}
        \hline
        \textbf{Tham số} & \textbf{Giá trị} \\
        \hline
        \texttt{event\_top\_k} & \(8\) \\
        \texttt{direct\_rule\_top\_k} & \(8\) \\
        \texttt{semantic\_pool\_size} & \(100\) \\
        \texttt{max\_hops} & \(2\) \\
        \texttt{max\_paths\_per\_event} & \(30\) \\
        \texttt{max\_candidate\_rules} & \(200\) \\
        \texttt{final\_top\_k} & \(12\) \\
        \texttt{min\_event\_score} & \(0{,}20\) \\
        \texttt{min\_rule\_score} & \(0{,}15\) \\
        \texttt{direction} & \texttt{both} \\
        \hline
    \end{tabular}
\end{table}

Tùy chọn \texttt{direction=both} cho phép retriever duyệt graph theo cả chiều
thuận và chiều ngược. Việc này cần thiết vì benchmark bao gồm cả forward
multihop và reverse reasoning. Tuy nhiên, việc cho phép hai hướng không bảo
đảm path được xếp hạng phù hợp với answer intent; ảnh hưởng này được phân
tích trong Phần~\ref{sec:discussion}.

\subsection{Cấu hình query-aware verification}

Các hyperparameter chung của Step 4 được trình bày trong
Bảng~\ref{tab:verification-hyperparameters}.

\begin{table}[t]
    \centering
    \caption{Hyperparameter của query-aware verification.}
    \label{tab:verification-hyperparameters}
    \begin{tabular}{lr}
        \hline
        \textbf{Tham số} & \textbf{Giá trị} \\
        \hline
        \texttt{cf\_top\_k} & \(5\) \\
        \texttt{mapping\_top\_k} & \(5\) \\
        \texttt{mapping\_threshold} & \(0{,}42\) \\
        \texttt{max\_cf\_hops} & \(3\) \\
        \texttt{max\_cf\_paths} & \(30\) \\
        \texttt{verified\_top\_k} & \(10\) \\
        \texttt{keep\_threshold} & \(0{,}52\) \\
        \texttt{reject\_threshold} & \(0{,}34\) \\
        Alternative path threshold & \(0{,}35\) \\
        Semantic mapping & Bật \\
        \hline
    \end{tabular}
\end{table}

Semantic mapping được bật trong cả hai run. Vì vậy, sự khác biệt giữa
structural SCM và path ablation không xuất phát từ việc một cấu hình sử dụng
embedding mapping còn cấu hình kia thì không.

\subsection{Hai chế độ xác minh}

Paired experiment gồm hai cấu hình.

\subsubsection{Structural SCM}

Cấu hình thứ nhất sử dụng:

\begin{verbatim}
counterfactual_mode = structural_scm
\end{verbatim}

LegalSCM thực hiện:

\begin{itemize}
    \item factual path validation;
    \item kiểm tra rule coverage;
    \item suy luận fixed point ba giá trị;
    \item hard intervention trên mediator;
    \item suy luận lại counterfactual outcome;
    \item lưu disabled rule, newly activated rule và intervention trace.
\end{itemize}

Node-deletion baseline vẫn được lưu trong trace dưới dạng diagnostic, nhưng
final structural result được ưu tiên khi LegalSCM xử lý thành công.

\subsubsection{Path ablation}

Cấu hình thứ hai sử dụng:

\begin{verbatim}
counterfactual_mode = path_ablation
\end{verbatim}

Mediator được xóa khỏi graph và hệ thống tìm bounded alternative path từ
seed event đến outcome. Cấu hình này không thực hiện structural fixed-point
inference và không được diễn giải như phép \(do(\cdot)\).

\subsubsection{Điều kiện paired comparison}

Hai run sử dụng cùng:

\begin{itemize}
    \item \(250\) câu hỏi benchmark;
    \item retriever và retrieval output schema;
    \item embedding model và semantic mapping;
    \item graph, rule corpus, memory và index;
    \item verification threshold;
    \item Step 5 answer generator;
    \item extractive provider;
    \item evaluator và metric.
\end{itemize}

Biến độc lập duy nhất là \texttt{counterfactual\_mode}. Vì vậy, chênh lệch
giữa hai run có thể được quy về verification engine trong phạm vi
implementation hiện tại.

\subsection{Cấu hình sinh câu trả lời}

Các hyperparameter của Step 5 được trình bày trong
Bảng~\ref{tab:generation-hyperparameters}.

\begin{table}[t]
    \centering
    \caption{Hyperparameter của answer generation.}
    \label{tab:generation-hyperparameters}
    \begin{tabular}{lr}
        \hline
        \textbf{Tham số} & \textbf{Giá trị} \\
        \hline
        Provider & \texttt{extractive} \\
        \texttt{max\_evidence} & \(8\) \\
        \texttt{answer\_max\_paths} & \(6\) \\
        \texttt{max\_context\_chars} & \(18\,000\) \\
        \texttt{max\_tokens} & \(1\,200\) \\
        \texttt{temperature} & \(0{,}1\) \\
        \texttt{timeout} & \(180\) giây \\
        \texttt{min\_verification\_score} & \(0{,}45\) \\
        Include uncertain evidence & Không \\
        Extractive fallback & Bật \\
        \hline
    \end{tabular}
\end{table}

Do provider chính là extractive, các tham số \texttt{max\_tokens} và
\texttt{temperature} không ảnh hưởng trực tiếp đến nội dung câu trả lời
trong paired experiment. Chúng được giữ trong configuration để cùng schema
có thể hỗ trợ các provider sinh khác.

Không sử dụng LLM judge hoặc LLM generation trong kết quả chính. Lựa chọn
này giúp giảm biến thiên ngẫu nhiên và cho phép phân tích trực tiếp ảnh
hưởng của retrieval, verification và answer intent.

\subsection{Cấu hình evaluator}

Step 7 được chạy với các giá trị:

\[
k \in \{1,3,5,10\}.
\]

Evaluator sử dụng:

\begin{itemize}
    \item toàn bộ \(250\) gold sample;
    \item toàn bộ \(250\) prediction;
    \item chế độ \texttt{strict};
    \item cùng một metric implementation cho hai cấu hình;
    \item thư mục output riêng cho structural SCM và path ablation.
\end{itemize}

Các thư mục kết quả chính là:

\begin{verbatim}
evaluation_results/structural_scm_v41_step5_v51/
evaluation_results/path_ablation_v41_step5_v51/
\end{verbatim}

Các report cũ không có hậu tố \texttt{step5\_v51} không được sử dụng trong
paired comparison.

\subsection{Metric retrieval}

Gọi \(\mathcal{G}\) là tập gold ID và
\(\widehat{\mathcal{R}}_k\) là \(k\) ID được xếp hạng cao nhất. Precision và
Recall tại \(k\) được tính bởi:

\begin{equation}
\mathrm{Precision@}k
=
\frac{
    \left|
        \widehat{\mathcal{R}}_k
        \cap
        \mathcal{G}
    \right|
}{
    k
},
\end{equation}

\begin{equation}
\mathrm{Recall@}k
=
\frac{
    \left|
        \widehat{\mathcal{R}}_k
        \cap
        \mathcal{G}
    \right|
}{
    \left|\mathcal{G}\right|
}.
\end{equation}

Hit@\(k\) nhận giá trị \(1\) nếu ít nhất một gold item xuất hiện trong top
\(k\), ngược lại nhận giá trị \(0\).

Các metric được tính riêng cho:

\begin{itemize}
    \item rule ID;
    \item event ID;
    \item article ID.
\end{itemize}

Evaluator còn báo cáo:

\begin{itemize}
    \item Macro Precision, Recall và F1;
    \item Micro Precision, Recall và F1;
    \item Exact Set Match;
    \item Mean Reciprocal Rank (MRR);
    \item Mean Average Precision (MAP).
\end{itemize}

Trong paper, Rule Recall@5 và Event Recall@5 được sử dụng làm các retrieval
metric chính vì benchmark có thể có nhiều gold rule hoặc event trên cùng một
path.

\subsection{Metric causal path}

Gold path và predicted path được chuẩn hóa thành chuỗi cạnh có hướng:

\[
\left[
(v_0,v_1),
(v_1,v_2),
\ldots,
(v_{h-1},v_h)
\right].
\]

\paragraph{Top-1 Exact Path Match.}

Metric nhận giá trị \(1\) nếu path được pipeline chọn làm primary path trùng
hoàn toàn với gold path về thứ tự cạnh, ngược lại nhận giá trị \(0\).

\paragraph{Oracle Exact Path Match.}

Evaluator duyệt các candidate path được retriever trả về và chọn path khớp
gold tốt nhất. Oracle metric chỉ dùng để chẩn đoán khả năng retrieval; nó
không phản ánh path mà hệ thống thực sự sử dụng khi sinh câu trả lời.

Khoảng cách:

\[
\mathrm{OracleExact}
-
\mathrm{Top1Exact}
\]

được sử dụng để nhận diện lỗi path ranking. Nếu cả oracle và top-1 đều sai,
lỗi có khả năng nằm ở path retrieval.

\paragraph{Edge-level metrics.}

Predicted edge và gold edge được xem như hai tập để tính precision, recall
và F1.

\paragraph{Event-level metrics.}

Tập event xuất hiện trên predicted path được so sánh với tập event trên gold
path.

\paragraph{Final Event Accuracy.}

Metric kiểm tra outcome event cuối của predicted path có trùng với gold final
event hay không.

\paragraph{Path Rule F1.}

Tập rule hỗ trợ predicted path được so sánh với gold rule của path.

Đối với branch và convergence, exact path tuyến tính không phản ánh đầy đủ
cấu trúc cần dự đoán. Vì vậy, kết quả path metric của các nhóm này được báo
cáo nhưng phải diễn giải cùng answer và citation metric.

\subsection{Metric xác minh}

Verification decision có ba nhãn:

\[
\mathcal{D}
=
\{
\texttt{SUPPORTED},
\texttt{REJECT\_DIRECT\_CLAIM},
\texttt{UNCERTAIN}
\}.
\]

Accuracy được tính bằng:

\begin{equation}
\mathrm{Accuracy}
=
\frac{
    \text{số prediction có decision đúng}
}{
    \text{tổng số prediction}
}.
\end{equation}

Do benchmark mất cân bằng nhãn, evaluator báo cáo thêm:

\begin{itemize}
    \item precision, recall và F1 theo từng lớp;
    \item Macro-F1 trên các lớp có gold support;
    \item Balanced Accuracy;
    \item confusion matrix.
\end{itemize}

Balanced Accuracy là trung bình recall của các lớp có mặt trong gold data:

\begin{equation}
\mathrm{BalancedAccuracy}
=
\frac{1}{|\mathcal{D}^{+}|}
\sum_{d \in \mathcal{D}^{+}}
\mathrm{Recall}_d,
\end{equation}

trong đó \(\mathcal{D}^{+}\) là tập các lớp có ít nhất một gold sample.

Counterfactual verification accuracy được báo cáo riêng trên các mẫu có
\texttt{requires\_counterfactual=true}. Tuy nhiên, như đã trình bày tại
Phần~\ref{sec:data}, subset này chủ yếu đánh giá direct-edge negative claims.

\subsection{Metric câu trả lời}

Trước khi đánh giá, predicted answer và gold answer được:

\begin{itemize}
    \item chuyển về chữ thường;
    \item chuẩn hóa khoảng trắng;
    \item loại dấu câu không cần thiết;
    \item tách thành token;
    \item giữ nội dung tiếng Việt.
\end{itemize}

Token Precision và Token Recall được tính từ số token giao nhau giữa
prediction và gold:

\begin{equation}
P_{\mathrm{token}}
=
\frac{
    |\widehat{T} \cap T|
}{
    |\widehat{T}|
},
\end{equation}

\begin{equation}
R_{\mathrm{token}}
=
\frac{
    |\widehat{T} \cap T|
}{
    |T|
},
\end{equation}

trong đó \(\widehat{T}\) và \(T\) lần lượt là multiset token của prediction
và gold answer.

Answer Token F1 được tính bởi:

\begin{equation}
F1_{\mathrm{token}}
=
\frac{
    2P_{\mathrm{token}}R_{\mathrm{token}}
}{
    P_{\mathrm{token}}+R_{\mathrm{token}}
}.
\end{equation}

Nếu cả precision và recall bằng \(0\), F1 được đặt bằng \(0\).

ROUGE-L F1 được tính dựa trên longest common subsequence giữa prediction và
gold answer \cite{lin2004rouge}. Metric này giữ lại một phần thông tin về thứ
tự token mà bag-of-token F1 không phản ánh.

Evaluator còn báo cáo Normalized Exact Match. Tuy nhiên, do câu trả lời
extractive thường bổ sung causal chain và citation, Exact Match không được
chọn làm metric chính.

\subsection{Metric citation}

Citation được chuẩn hóa về article ID. Ví dụ:

\[
\text{``Điều 15''}
\longrightarrow
\texttt{15}.
\]

Các citation trùng lặp bị loại bỏ trước khi đánh giá. Với tập predicted
citation \(\widehat{\mathcal{C}}\) và gold citation \(\mathcal{C}\), evaluator
tính:

\begin{equation}
P_{\mathrm{citation}}
=
\frac{
    |
    \widehat{\mathcal{C}}
    \cap
    \mathcal{C}
    |
}{
    |
    \widehat{\mathcal{C}}
    |
},
\end{equation}

\begin{equation}
R_{\mathrm{citation}}
=
\frac{
    |
    \widehat{\mathcal{C}}
    \cap
    \mathcal{C}
    |
}{
    |
    \mathcal{C}
    |
},
\end{equation}

và harmonic F1 tương ứng. Citation Exact Match nhận giá trị \(1\) khi hai
tập article ID giống hoàn toàn.

Citation F1 đánh giá việc chọn đúng điều luật, nhưng không tự động chứng minh
rằng phần diễn giải của câu trả lời phù hợp với toàn bộ nội dung điều luật.

\subsection{Đo thời gian xử lý}

Batch runner lưu thời gian end-to-end của từng mẫu, bao gồm retrieval,
verification và answer generation. Step 7 báo cáo:

\begin{itemize}
    \item average runtime trên mỗi mẫu;
    \item median runtime trên mỗi mẫu.
\end{itemize}

Run log còn lưu wall-clock time của toàn bộ batch. Hai loại runtime được báo
cáo riêng:

\begin{itemize}
    \item per-sample runtime phản ánh thời gian xử lý được ghi cho từng câu;
    \item batch wall-clock time còn bao gồm khởi tạo model, đọc resource,
    checkpoint và ghi artifact.
\end{itemize}

Do thông tin phần cứng chưa được điền đầy đủ, runtime chỉ được sử dụng để so
sánh tương đối giữa hai cấu hình chạy trên cùng server, không được xem là
benchmark tốc độ có thể tổng quát sang môi trường khác.

\subsection{Phân tích lỗi}

Step 7 gán một hoặc nhiều error type cho từng sample. Các nhóm lỗi chính
gồm:

\begin{description}
    \item[\texttt{RULE\_RETRIEVAL\_MISS}]
    Gold rule không xuất hiện đầy đủ trong top-\(k\).

    \item[\texttt{EVENT\_RETRIEVAL\_MISS}]
    Gold event không được truy hồi đầy đủ.

    \item[\texttt{TOP1\_PATH\_MISMATCH}]
    Primary path không trùng gold path.

    \item[\texttt{PATH\_RANKING\_ERROR}]
    Gold path xuất hiện trong candidate set nhưng không được xếp đầu.

    \item[\texttt{PATH\_RETRIEVAL\_ERROR}]
    Gold path không xuất hiện trong candidate set.

    \item[\texttt{VERIFICATION\_ERROR}]
    Predicted decision khác gold decision.

    \item[\texttt{LOW\_ANSWER\_TOKEN\_F1}]
    Token overlap của câu trả lời thấp hơn ngưỡng phân tích lỗi.

    \item[\texttt{CITATION\_MISS}]
    Thiếu ít nhất một gold citation.

    \item[\texttt{EXTRA\_CITATION}]
    Câu trả lời chứa citation không thuộc gold set.

    \item[\texttt{PIPELINE\_ERROR}]
    Sample không tạo được prediction hoàn chỉnh.
\end{description}

Error type chỉ phục vụ phân tích sau inference và không được sử dụng để điều
chỉnh prediction.

\subsection{Phân tích theo nhóm}

Metric được tổng hợp theo ba cách:

\begin{enumerate}
    \item theo question type;
    \item theo difficulty;
    \item theo cờ \texttt{requires\_counterfactual}.
\end{enumerate}

Phân tích theo question type giúp xác định hệ thống mạnh hoặc yếu ở forward,
reverse, bridge, chain, branch và convergence. Phân tích theo difficulty
kiểm tra mức suy giảm trên các mẫu có nhiều candidate rule hoặc cấu trúc
phức tạp. Phân tích counterfactual subset cho phép kiểm tra riêng các direct
claim bị bác bỏ.

\subsection{Nguyên tắc diễn giải paired ablation}

Paired experiment được diễn giải theo các nguyên tắc sau:

\begin{enumerate}
    \item Nếu retrieval và path metric giống nhau, điều này được kỳ vọng vì
    verification engine được áp dụng sau retrieval.

    \item Nếu verification hoặc answer metric khác nhau, chênh lệch có thể
    liên quan đến final decision, primary path hoặc evidence được Step 4 giữ
    lại.

    \item Nếu quality metric giống nhau nhưng runtime khác nhau, không được
    kết luận structural SCM cải thiện độ chính xác.

    \item Trace chi tiết hơn là một lợi ích về auditability, không phải một
    quality improvement trừ khi được đánh giá bằng metric hoặc expert study
    riêng.

    \item Không sử dụng Oracle Exact Path thay cho Top-1 Exact Path khi báo
    cáo hiệu năng deployable.

    \item Không xem counterfactual accuracy trên \(20\) direct-edge negative
    samples là xác nhận đầy đủ cho hard-intervention reasoning.
\end{enumerate}

Do hai cấu hình tất định và chỉ được chạy một lần trên cùng benchmark, nghiên
cứu báo cáo descriptive paired comparison. Không thực hiện kiểm định ý nghĩa
thống kê hoặc báo cáo độ lệch chuẩn giữa nhiều random seed. Khi hai cấu hình
cho kết quả giống hệt nhau, kết quả được mô tả là \emph{null quality
ablation}, không phải bằng chứng về sự tương đương tổng quát trên mọi bộ dữ
liệu.
